\section{Fitting}
\label{sec:fitting}

\subsection{A population-robust objective}

The quantity we optimise is not the mean win rate against a fixed pool but the
worst win rate across a panel of opponents, softened so that a noisy search can
make progress on it. For a parameter vector $w$ and an opponent panel
$\mathcal{O}$, writing $\mathrm{wr}_o(w)$ for the win rate of the configuration
$w$ against opponent $o$ on the current seed bank, the fitting objective is
\begin{equation}
  \mathrm{score}(w) \;=\; -\frac{1}{\beta}\,
     \log \sum_{o \in \mathcal{O}} \exp\bigl(-\beta\, \mathrm{wr}_o(w)\bigr).
  \label{eq:softmin}
\end{equation}
This tends to $\min_o \mathrm{wr}_o(w)$ as $\beta$ grows and is smooth in $w$ for
finite $\beta$. The selection reported here used $\beta = 8$, the default of
\filepath{engine/select_final.py}. Two other values appear in the repository and
are not the reported ones: the tuner command line defaults to $\beta = 10$, and
the local response probe of \S\ref{sec:lbr} used $\beta = 1$, because that search
targets a single frozen opponent and the soft minimum degenerates to a mean.

The panel $\mathcal{O}$ used for fitting holds six policies: \vthree{}, the
turn-starvation lockout, the posterior detective, \vtwo{}, the adaptive
diversifier and the focused hunter. The misdirection artist and the random legal
control were excluded from every fitting and selection run so that they remain
held-out opponent identities. Two of the eight opponents in the evaluation panel
of \S\ref{sec:results} are therefore held-out identities and six are not, so the
head-to-head results measure generalisation to new deals rather than to new
strategic populations.

\subsection{Optimiser}

We use the cross-entropy method with a diagonal Gaussian over the parameter
vector. Each generation samples \numCEMPop{} candidates around the incumbent
mean, retains the \numCEMElite{} best as the elite set, and smooths the mean and
the per-coordinate standard deviation towards the elite statistics. The objective
is a noisy win rate estimated from a finite bank of deals, a regime in which
optimisers that trust small differences between candidates behave poorly. Every
candidate within a generation is evaluated on \emph{identical} seed banks, so
within-generation comparisons are paired and the deal draw cancels, and the banks
are redrawn between generations, so the search cannot lock onto a particular set
of deals.

The parameter vector is fitted jointly and has \numParams{} coordinates: the
\numAskFeats{} ask weights $w_0,\dots,w_{19}$ of Table~\ref{tab:features},
followed by \numKnobs{} decision knobs --- five thresholds and pool sizes
governing declaration and patience, the minimum team probability, the two mixing
weights $\lambda_{\mathrm{val}}$ and $\lambda_{\mathrm{lin}}$ of
\eqref{eq:askutil}, the stopping margin $\varepsilon$ of \eqref{eq:stop}, the two
policy-prior coefficients $\theta$ and $\phi$ of \eqref{eq:prior}, and the
lookahead width $K$ with the chain and threat weights of \S\ref{sec:twoply}.
Bounds are imposed per coordinate so that probabilities stay in $[0,1]$, counts
stay non-negative integers, and the lookahead width stays within its
implementation limit; every knob is named with its bound in
Table~\ref{tab:bounds} of \S\ref{app:params}.

The value function of \S\ref{sec:value} is not part of this vector. It was fitted
separately by ridge regression on self-play decision points and then held fixed
while the policy weights were fitted around it, so the search never adjusted the
value coefficients and the policy weights are conditional on them.

\subsection{Schedule and selection}

Fitting ran in five rounds. Round one fitted the \numAskFeats{} ask weights alone
against the Fast posterior on a four-opponent panel, from base seed $20260821$;
round two started from its mean and fitted the joint vector on the six-opponent
panel, from base seed $770077$, and is superseded because its coordinate layout
predates the two ask features added afterwards; round three re-ran the joint fit
on the aligned \numParams{}-coordinate vector from base seed $313131$; and round
four repeated it from base seed $888111$ after a correction to the
information-leak features. Round five produced the shipping configuration. Its
base seed is not recorded in any committed script or artifact; this is a
reproducibility gap, and the round can be inspected but not reproduced exactly
from the trace \filepath{research/v04/runs/tune-round5.jsonl}.

The shipping configuration was chosen on a validation bank the search itself
never saw, because the per-generation best reported by the optimiser is a maximum
over a population evaluated on shared seeds and is therefore biased upwards. The
last six generation means of the round-five trace were re-evaluated on seed
$1357911$ (500 deals, two rotations, panel \{\vthree{}, lockout, detective,
\vtwo{}\}), and generation \numSelectedGen{} won that selection, with a
validation soft-min of \numSelectSoftmin{} and per-opponent win rates of
\numSelectPerOpp{}, as recorded in \filepath{research/v04/runs/selected.json}.
Every seed bank used from \S\ref{sec:protocol} onwards is disjoint from the
fitting and selection banks, and the configuration was frozen before any of those
runs.

Table~\ref{tab:params} in \S\ref{app:params} lists the frozen values. It is a
reproducibility artifact and should not be read as a description of strategy: the
ask features are correlated with one another and are normalised to different
effective ranges, so the sign and magnitude of a single coefficient do not
identify the contribution of the corresponding feature. Two of the decision
knobs --- the locked-half-suit threshold and the patience switch --- are inert
while the value-based rule of \eqref{eq:stop} is active, so their fitted values
carry no behavioural meaning in the shipping configuration.
